Gemini models build on top of Transformer decoders \citep{vaswani2017} that are enhanced with improvements in architecture and model optimization to enable stable training at scale and optimized inference on Google's Tensor Processing Units. They are trained to support 32k context length, employing efficient attention mechanisms (for e.g. multi-query attention \citep{mqa}). Our first version, Gemini 1.0, comprises three main sizes to support a wide range of applications as discussed in Table~\ref{tab:family}.

\begin{table}[h!]
\begin{tabularx}{\textwidth}{lX}
\toprule
\textbf{Model size} & \textbf{Model description} \\
\midrule
Ultra & Our most capable model that delivers state-of-the-art performance across a wide range of highly complex tasks, including reasoning and multimodal tasks. It is efficiently serveable at scale on TPU accelerators due to the Gemini architecture. \\
\midrule
Pro & A performance-optimized model in terms of cost as well as latency that delivers significant performance across a wide range of tasks. This model exhibits strong reasoning performance and broad multimodal capabilities. \\
\midrule
Nano & Our most efficient model, designed to run on-device. We trained two versions of Nano, with 1.8B (Nano-1) and 3.25B (Nano-2) parameters,  targeting low and high memory devices respectively. It is trained by distilling from larger Gemini models. It is 4-bit quantized for deployment and provides best-in-class performance.  \\
\bottomrule
%\end{tabular}
\end{tabularx}
 \caption{An overview of the Gemini 1.0 model family.}
    \label{tab:family}
\end{table}

Gemini models are trained to accommodate textual input interleaved with a wide variety of audio and visual inputs, such as natural images, charts, screenshots, PDFs, and videos, and they can produce text and image outputs (see Figure~\ref{fig:arch}). The visual encoding of Gemini models is inspired by our own foundational work on Flamingo \citep{flamingo}, CoCa \cite{coca}, and PaLI \citep{pali}, with the important distinction that the models are multimodal from the beginning and can natively output images using discrete image tokens \cite{dalle, parti}.

\begin{figure}[h!]
\centering
 \includegraphics[width=0.9\textwidth]{figs/tech_fig_architecture.png}
 \caption{Gemini models support interleaved sequences of text, image, audio, and video as inputs (illustrated by tokens of different colors in the input sequence). They can output responses with interleaved image and text.}
 \label{fig:arch}
\end{figure}

Video understanding is accomplished by encoding the video as a sequence of frames in the large context window. Video frames or images can be interleaved naturally with text or audio as part of the model input. The models can handle variable input resolution in order to spend more compute on tasks that require fine-grained understanding. In addition, Gemini models can directly ingest audio signals at 16kHz from Universal Speech Model (USM)~\citep{usm} features. This enables the model to capture nuances that are typically lost when the audio is naively mapped to a text input (for example, see audio understanding demo on the \href{https://deepmind.google/gemini
}{website}).

Training the Gemini family of models required innovations in training algorithms, dataset, and infrastructure. For the Pro model, the inherent scalability of our infrastructure and learning algorithms enable us to complete pre-training in a matter of weeks, leveraging a fraction of the Ultra's resources. The Nano series of models leverage additional advancements in distillation and training algorithms to produce the best-in-class small language models for a wide variety of tasks, such as summarization and reading comprehension, which power our next generation on-device experiences.